% =====================================================================
%  CHAPTER 13: 社会卫生状况
% =====================================================================
\chapter{社会卫生状况}
\setcounter{chapter}{13}
\chapterintro{
掌握人群健康状况的复合型评价指标（PYLL、DALY等）和社会发展指标（HDI、ASHA、PQLI等）（\textbf{教学难点}）；掌握初级卫生保健的概念、基本原则和要素；熟悉社会卫生状况的概念和评价程序；了解中国及全球社会卫生状况、全球卫生策略和可持续发展目标。
}

% ----- 13.1 社会卫生状况概述 -----
\section{社会卫生状况概述}

\begin{defbox}
\textbf{社会卫生状况（Social Health Status）}：一定时空范围内，人群的健康状况以及影响人群健康的各种社会因素的综合状况。

\textbf{人群健康状况}：人群的整体健康水平，通过死亡率、发病率、期望寿命等一系列指标来测量和评价。

\textbf{影响人群健康的社会因素：}
\begin{itemize}[leftmargin=*]
    \item 社会经济发展水平
    \item 人口特征（数量、结构、素质、流动）
    \item 社会文化背景（教育、文化、宗教、风俗）
    \item 生活和工作环境条件
    \item 卫生服务系统（资源配置、可及性、质量）
\end{itemize}
\end{defbox}

\subsection{社会卫生状况评价}

\begin{enumerate}[leftmargin=*]
    \item \imp{评价资料来源}：①人口普查数据；②生命登记系统（出生、死亡统计）；③疾病监测系统；④卫生服务常规报表；⑤专题调查（如国家卫生服务调查）；⑥国际组织数据库（WHO、UNICEF、世界银行等）
    \item \imp{评价程序}：确定评价目标→选择评价指标→收集数据→分析与比较（横向对比不同地区、纵向比较不同时期）→发现问题与提出建议
    \item \imp{评价意义}：①全面了解人群健康状况和主要卫生问题；②明确卫生工作的重点和优先领域；③评价卫生政策和干预措施的效果；④为卫生资源配置提供依据
\end{enumerate}

% ----- 13.2 人群健康状况评价指标 -----
\section{人群健康状况的评价指标}

\subsection{单一型指标}

\begin{table}[htbp]
\centering
\caption{常用单一型健康指标}
\begin{tabularx}{\textwidth}{lX}
\hline
\textbf{指标类型} & \textbf{具体指标} \\
\hline
\imp{死亡指标} & 死亡率、年龄别死亡率、婴儿死亡率（IMR）、5岁以下儿童死亡率（U5MR）、孕产妇死亡率（MMR）、死因别死亡率 \\
\imp{发病指标} & 发病率、患病率、疾病构成比 \\
\imp{残疾指标} & 残疾发生率、残疾现患率 \\
\imp{营养指标} & 低出生体重率、5岁以下儿童低体重率、发育迟缓率 \\
\imp{期望寿命} & 出生期望寿命（Life Expectancy at Birth）、健康期望寿命（HALE） \\
\hline
\end{tabularx}
\end{table}

\subsection{复合型指标（教学难点*）}

\begin{keybox}
\imp{（教学难点*）}\textbf{主要复合型健康指标：}

\textbf{1. 潜在寿命损失年（Potential Years of Life Lost, PYLL）：}
\begin{itemize}[leftmargin=*]
    \item 定义：某年龄组人群因某病死亡而损失的寿命年数
    \item 计算公式：PYLL = $\sum$ D\textsubscript{x} $\times$ (L $-$ x)，其中D\textsubscript{x}=年龄x的死亡人数，L=标准寿命上限（通常为70岁），x=死亡年龄的中点
    \item 意义：强调\imp{"早死"}对寿命的损失，给予年轻死亡更高的权重
    \item 应用：确定优先干预的疾病和健康问题
\end{itemize}

\textbf{2. 无残疾期望寿命（Disability-Free Life Expectancy, DFLE/LEFD）：}
\begin{itemize}[leftmargin=*]
    \item 以无残疾状态下的期望寿命来评价人群健康水平
    \item 扣除了因残疾而损失的寿命年
\end{itemize}

\textbf{3. 活动期望寿命（Active Life Expectancy, ALE）：}
\begin{itemize}[leftmargin=*]
    \item 日常生活中能保持自理能力、无依赖的期望寿命年限
    \item 对老年人群的健康评价尤为重要
\end{itemize}

\textbf{4. 伤残调整寿命年（Disability Adjusted Life Year, DALY）：}
\begin{itemize}[leftmargin=*]
    \item \keyword{目前最广泛使用的疾病负担综合评价指标}
    \item DALY = YLLs（早死生命损失年）+ YLDs（伤残生命损失年）
    \item 综合考虑了早死和失能两种健康损失
    \item 用于疾病负担（Global Burden of Disease, GBD）研究
    \item 意义：一个DALY代表因早死或伤残而损失的一年"完全健康"生命
    \item 可比较不同疾病、不同地区、不同时期的疾病负担
    \item \textbf{DALY的局限性：}①年龄权重和贴现率的选择有争议；②残疾权重的主观性；③未考虑共存疾病的影响；④数据不足影响低收入国家的精确性；⑤未包括社会环境因素对疾病负担的影响
\end{itemize}

\textbf{5. 健康期望寿命（Healthy Life Expectancy, HALE）：}
\begin{itemize}[leftmargin=*]
    \item 以"完全健康状态"的期望寿命来衡量人群健康水平
    \item WHO自2000年起定期发布各成员国的HALE数据
\end{itemize}
\end{keybox}

% ----- 13.3 健康影响因素指标 -----
\section{健康影响因素指标}

\subsection{人口学指标}

\begin{itemize}[leftmargin=*]
    \item \imp{人口数量}：总人口数、人口密度
    \item \imp{人口分布}：城乡分布、地区分布
    \item \imp{性别构成}：性别比（正常103-107）
    \item \imp{年龄构成}：各年龄组占比、老龄化程度（65岁及以上$\geq$7\%为老龄化）、抚养系数
    \item \imp{人口自然增长}：出生率-死亡率=自然增长率
\end{itemize}

\subsection{社会发展指标（教学难点*）}

\begin{keybox}
\imp{（教学难点*）}\textbf{主要社会发展指标：}

\textbf{1. 人类发展指数（Human Development Index, HDI）：}
\begin{itemize}[leftmargin=*]
    \item UNDP于1990年提出，综合反映社会经济发展水平
    \item 三个维度：健康长寿（预期寿命指数）、教育获得（教育指数）、生活水平（收入指数，按购买力平价调整的人均GDP）
    \item HDI = 三个维度指数的几何平均数
    \item 分级：$\geq$0.800为极高人类发展水平
\end{itemize}

\textbf{2. ASHA指数（American Social Health Association Index）：}
\begin{itemize}[leftmargin=*]
    \item 由美国社会健康协会提出
    \item 计算公式：ASHA = （就业率 $\times$ 识字率 $\times$ 平均期望寿命 $\times$ 人均GDP增长率）/ （出生率 $\times$ 婴儿死亡率）
    \item 综合反映社会经济发展和健康水平
\end{itemize}

\textbf{3. 物质生活质量指数（Physical Quality of Life Index, PQLI）：}
\begin{itemize}[leftmargin=*]
    \item 由婴儿死亡率、1岁期望寿命和识字率三项指标组成
    \item 三项指标各赋0-100分，求算术平均值
    \item 值越高表明生活质量越好
    \item 优势：简单、易于计算、不受GDP影响
\end{itemize}

\textbf{4. 社会进步指数（Social Progress Index, SPI）：}
\begin{itemize}[leftmargin=*]
    \item 衡量社会在满足人类基本需求、建立福祉基础和创造机会三个方面表现的综合性指标
\end{itemize}
\end{keybox}

\subsection{健康公平指标}

\begin{defbox}
\textbf{健康公平（Health Equity）}：不同社会群体之间不存在可避免的、不公平的健康差异。健康公平关注的是\keyword{健康差异的社会根源和可改变性}。

\textbf{主要衡量指标：}
\begin{itemize}[leftmargin=*]
    \item \imp{基尼系数（Gini Coefficient）}：取值范围0-1，衡量健康或卫生服务在人群中的不平等程度。0=完全平等，1=完全不平等
    \item \imp{集中指数（Concentration Index）}：衡量与经济社会地位相关的健康不平等程度。正值表示健康集中在富人群体，负值表示集中在穷人群体
    \item \imp{比率和差值}：不同阶层或地区健康指标之比和之差，直观反映差异大小
\end{itemize}
\end{defbox}

% ----- 13.4 中国与全球社会卫生状况 -----
\section{中国与全球社会卫生状况}

\subsection{中国社会卫生状况}

\begin{keybox}
\textbf{健康中国（Healthy China）——中国卫生事业发展战略：}

\begin{enumerate}[leftmargin=*]
    \item \textbf{《"健康中国2030"规划纲要》（2016年）}——将健康纳入国家发展战略的核心，提出"以人民健康为中心"的发展理念
    \item \textbf{五大战略任务：}普及健康生活、优化健康服务、完善健康保障、建设健康环境、发展健康产业
    \item \textbf{核心目标：}2030年人均期望寿命达到79岁（2020年77.3岁），婴儿死亡率降至5.0‰以下，孕产妇死亡率降至12/10万以下
    \item \textbf{主要成就：}人均期望寿命从1949年的35岁升至2020年的77.3岁；孕产妇死亡率从1500/10万降至16.9/10万；婴儿死亡率从200‰降至5.4‰
\end{enumerate}
\end{keybox}

\subsection{全球健康状况与全球卫生}

\begin{itemize}[leftmargin=*]
    \item 全球健康面临的核心挑战：传染病（艾滋病、结核病、疟疾）+ 慢性非传染性疾病（NCDs占全球死亡71\%）+ 新发传染病 + 健康不公平
    \item 全球卫生：超越国界的、需要全球共同行动来解决的健康问题和挑战
\end{itemize}

% ----- 13.5 全球卫生策略 -----
\section{全球卫生策略}

\subsection{21世纪人人享有卫生保健}

\begin{keybox}
\textbf{WHO"21世纪人人享有卫生保健"（HFA 21世纪）：}

\textbf{主要内容：}
\begin{enumerate}[leftmargin=*]
    \item 实现健康公平——缩小国家间和国家内部的健康差距
    \item 强调生命全程健康——从围生期到老年的全程保障
    \item 预防为主——重视非传染性疾病和精神卫生的预防
    \item 改善卫生系统绩效——提高质量、效率、反应性和公平性
\end{enumerate}

\textbf{总目标（至2020年）：}
\begin{enumerate}[leftmargin=*]
    \item 增加全体人民的期望寿命和生活质量
    \item 改进国家间和国家内部的健康公平
    \item 使全体人民能利用可持续的卫生系统服务
\end{enumerate}
\end{keybox}

\subsection{初级卫生保健（PHC）}

\begin{defbox}
\textbf{初级卫生保健（Primary Health Care, PHC）}：\keyword{1978年Alma-Ata宣言}正式提出，是基本的卫生保健服务，是社区中个人和家庭能够普遍获得的基本卫生服务，是卫生系统的最前沿和第一级接触点，也是实现"人人享有卫生保健"目标的关键途径。
\end{defbox}

\begin{keybox}
\textbf{初级卫生保健的基本原则：}
\begin{enumerate}[leftmargin=*]
    \item \imp{社会公平}——卫生服务应公平可及，特别关注弱势群体
    \item \imp{社区参与}——社区和个人应参与卫生服务的规划和组织
    \item \imp{成本效益}——采用技术上适宜、经济上可承担的方法和技术
    \item \imp{部门协作}——卫生部门之外的农业、教育、住房、公共工程等部门均应参与
    \item \imp{预防为主}——关注健康促进和疾病预防，而非仅仅是治疗
\end{enumerate}

\textbf{初级卫生保健的基本要素（8项）：}
\begin{enumerate}[leftmargin=*]
    \item 健康教育——普及健康知识
    \item 合理营养与食品安全——促进适当的食物供应和营养
    \item 安全饮水和基本环境卫生——足够的安全用水和基本卫生设施
    \item 妇幼保健与计划生育——包括围产期保健
    \item 计划免疫——主要传染病的免疫接种
    \item 地方病的预防和控制
    \item 常见病和伤害的适当治疗
    \item 基本药物的提供
\end{enumerate}
\end{keybox}

\subsection{千年发展目标与可持续发展目标}

\begin{table}[htbp]
\centering
\caption{千年发展目标（MDGs, 2000-2015）与可持续发展目标（SDGs, 2016-2030）的比较}
\begin{tabularx}{\textwidth}{lX}
\hline
\textbf{MDGs（8项目标）} & \textbf{SDGs（17项目标）} \\
\hline
消除极端贫困和饥饿；普及初等教育；促进性别平等；降低儿童死亡率；改善孕产妇健康；抗击艾滋病、疟疾等；确保环境可持续性；全球发展伙伴关系 &
目标3直接关注健康——"确保健康的生活方式，促进各年龄段所有人的福祉"；其他目标也影响健康——消除贫困（目标1）、零饥饿（目标2）、清洁饮水和卫生设施（目标6）、减少不平等（目标10）等；目标更广泛、更具普遍性 \\
\hline
\end{tabularx}
\end{table}

\begin{tipbox}
\textbf{可持续发展目标（SDGs）中与健康相关的目标：}\\
SDG 3 的具体目标包括——降低全球孕产妇死亡率（$<$70/10万）；消除新生儿和5岁以下儿童可预防的死亡；终结艾滋病、结核病、疟疾等传染病流行；将NCDs过早死亡减少1/3；实现全民健康覆盖（UHC）；确保普遍获得性与生殖健康服务。
\end{tipbox}

% ----- 复习题 -----
\section{复习题}

\subsection*{一、选择题}
\begin{enumerate}[leftmargin=*, itemsep=3pt]
    \item DALY等于：
    \begin{enumerate}
        \item[(A)] YLLs $-$ YLDs
        \item[(B)] YLLs $+$ YLDs
        \item[(C)] YLLs $\times$ YLDs
        \item[(D)] YLLs $\div$ YLDs
    \end{enumerate}
    \item PQLI由哪三项指标构成？
    \begin{enumerate}
        \item[(A)] 婴儿死亡率、1岁期望寿命、识字率
        \item[(B)] 婴儿死亡率、成人死亡率、GDP
        \item[(C)] 识字率、GDP、期望寿命
        \item[(D)] GDP、IMR、MMR
    \end{enumerate}
    \item 初级卫生保健（PHC）是哪次会议正式提出的？
    \begin{enumerate}
        \item[(A)] 1948年WHO宪章
        \item[(B)] 1978年Alma-Ata宣言
        \item[(C)] 2000年联合国千年首脑会议
        \item[(D)] 2015年联合国可持续发展峰会
    \end{enumerate}
    \item 中国政府提出的"健康中国2030"核心目标中，2030年人均期望寿命目标为：
    \begin{enumerate}
        \item[(A)] 75岁
        \item[(B)] 77岁
        \item[(C)] 79岁
        \item[(D)] 82岁
    \end{enumerate}
    \item SDGs中含多少项目标？
    \begin{enumerate}
        \item[(A)] 8项
        \item[(B)] 12项
        \item[(C)] 17项
        \item[(D)] 21项
    \end{enumerate}
\end{enumerate}

\subsection*{二、名词解释}
\begin{enumerate}[leftmargin=*, itemsep=3pt]
    \item \textbf{伤残调整寿命年（DALY）}
    \item \textbf{潜在寿命损失年（PYLL）}
    \item \textbf{初级卫生保健（PHC）}
    \item \textbf{健康公平（Health Equity）}
    \item \textbf{人类发展指数（HDI）}
\end{enumerate}

\subsection*{三、简答题}
\begin{enumerate}[leftmargin=*, itemsep=3pt]
    \item 简述DALY的构成及其用途。
    \item 简述初级卫生保健的五项基本原则和八项基本要素。
    \item 比较MDGs与SDGs在健康相关目标上的主要区别。
    \item 简述"健康中国2030"战略的五大战略任务。
\end{enumerate}

\subsection*{参考答案要点}

\subsubsection*{一、选择题}
1. (B)；2. (A)；3. (B)；4. (C)；5. (C)

\subsubsection*{二、名词解释}
\begin{enumerate}[leftmargin=*]
    \item \textbf{DALY}：伤残调整寿命年，综合反映早死（YLLs）和伤残（YLDs）造成的健康损失，是目前使用最广泛的疾病负担综合评价指标。
    \item \textbf{PYLL}：潜在寿命损失年，强调"早死"造成的寿命损失，给予年轻死亡更高的权重，用于确定优先干预领域。
    \item \textbf{PHC}：1978年Alma-Ata宣言提出的基本卫生保健，是社区中个人和家庭能普遍获得的基本卫生服务，是实现"人人享有卫生保健"的关键。
    \item \textbf{健康公平}：不同社会群体之间不存在可避免的、不公平的健康差异，关注健康差异的社会根源和可改变性。
    \item \textbf{HDI}：UNDP 1990年提出，综合反映社会经济发展水平的指标，包括预期寿命、教育获得和生活水平三个维度。
\end{enumerate}

\subsubsection*{三、简答题}
\begin{enumerate}[leftmargin=*]
    \item DALY = YLLs（早死生命损失年）+ YLDs（伤残生命损失年）。用途：比较不同疾病/地区/时期的疾病负担；为卫生资源配置和优先干预决策提供依据；监测卫生策略效果。局限性：年龄权重选择争议、残疾权重主观性、数据不足问题等。
    \item 五项原则：社会公平、社区参与、成本效益、部门协作、预防为主。八项要素：健康教育、合理营养与食品安全、安全饮水和环境卫生、妇幼保健与计划生育、计划免疫、地方病防控、常见病和伤害适当治疗、基本药物提供。
    \item MDGs（8项目标）主要针对发展中国家，聚焦于贫困、疾病等基本问题；SDGs（17项目标）面向所有国家，更具普遍性和综合性，健康目标（SDG3）涵盖面更广，包括NCDs、全民健康覆盖、精神卫生等。
    \item 五大任务：普及健康生活（健康教育、健康促进）、优化健康服务（强化公共卫生和基本医疗服务）、完善健康保障（健全医疗保障体系）、建设健康环境（环境保护、食品安全）、发展健康产业（医药、健康服务业）。
\end{enumerate}

\newpage
